\chapter{Practice workshops and worked solutions}
\label{learn:workshops}
\learning{Turn reference knowledge into repeatable tasks: query a business model, change data safely, use parameters, inspect a plan and prove that a backup can be restored. Each workshop includes a result or acceptance condition.}

\section{How to use the workshops}
Start with a fresh teaching database loaded by
\path{docs/manual/examples/01-create-shop.sql}. For query-only workshops the dataset
should remain unchanged. For experiments with schema or administration,
use a separate disposable database. Read the task before looking at its
solution. A correct answer must both produce the intended result and respect
the current implementation's limits.

These workshops deliberately distinguish correctness from speed. A four-row
table is suitable for proving join semantics; it is not suitable for drawing
conclusions about throughput. Likewise, a backup command's success status
proves less than an actual restore followed by a checker run and application
queries.

\section{Workshop 1: a stock and customer report}
\subsection{Tasks}
List all active customers in name order. Identify customers whose city is
unknown. Find the two products with the highest stock. Count the total
stock across all products. Explain why an empty city string and SQL NULL
would produce different answers to \code{COUNT(city)}.

\subsection{Solution}
\begin{lstlisting}[language=SQL]
SELECT customer_id, name
FROM customers WHERE active = TRUE ORDER BY name;

SELECT customer_id, name
FROM customers WHERE city IS NULL ORDER BY customer_id;

SELECT product_id, name, stock
FROM products ORDER BY stock DESC, product_id LIMIT 2;

SELECT SUM(stock) AS total_stock FROM products;
\end{lstlisting}
The active customers are Ada, Grace and Linus. Margaret has an unknown city.
The highest-stock products are USB-C Cable (100) and Notebook (50). Total
stock is 200. \code{COUNT(city)} ignores SQL NULL but counts an empty string
as a non-null value. Unknown and known-empty are different facts.

\subsection{Reasoning check}
A result may coincidentally appear sorted without an \code{ORDER BY}.
That does not make the ordering part of the query contract. If two products
have equal stock, the product identifier in the sort makes the displayed
order repeatable.

\section{Workshop 2: joins and customers without orders}
\subsection{Tasks}
Produce one row per order containing the customer's name. Then produce one
row per customer with an order count, including zero. Finally list only
customers with no orders. Explain why a left join followed by
\code{WHERE o.order\_id IS NULL} answers the last question.

\subsection{Solution}
\begin{lstlisting}[language=SQL]
SELECT o.order_id, c.name
FROM orders o JOIN customers c
  ON c.customer_id = o.customer_id
ORDER BY o.order_id;

SELECT c.customer_id, c.name, COUNT(o.order_id) AS orders
FROM customers c LEFT JOIN orders o
  ON o.customer_id = c.customer_id
GROUP BY c.customer_id, c.name
ORDER BY c.customer_id;

SELECT c.customer_id, c.name
FROM customers c LEFT JOIN orders o
  ON o.customer_id = c.customer_id
WHERE o.order_id IS NULL
ORDER BY c.customer_id;
\end{lstlisting}
The counts are 2, 1, 1 and 0 for customer identifiers 1--4. The final query
returns only Margaret. A matching order's primary key cannot be null.
Therefore the null right-hand key identifies the synthetic unmatched row
created by the outer join. Using a nullable non-key column as the unmatched
test would be unsafe because a real matched row might also contain null.

\section{Workshop 3: aggregates, CTEs and windows}
\subsection{Tasks}
Find the quantity in each order. Name the intermediate aggregate with a CTE,
then join it to orders and customers. Number each customer's orders in order
identifier order. Explain the difference between aggregation and a window
function in the shape of the result.

\subsection{Solution}
\begin{lstlisting}[language=SQL]
WITH units AS (
  SELECT order_id, SUM(quantity) AS unit_count
  FROM order_items GROUP BY order_id
)
SELECT o.order_id, c.name, u.unit_count
FROM orders o
JOIN customers c ON c.customer_id = o.customer_id
JOIN units u ON u.order_id = o.order_id
ORDER BY o.order_id;

SELECT order_id, customer_id,
  ROW_NUMBER() OVER (
    PARTITION BY customer_id ORDER BY order_id
  ) AS customer_order
FROM orders ORDER BY customer_id, order_id;
\end{lstlisting}
The first query returns order quantities 3, 3, 3 and 1. The window query
returns four order rows: Ada's orders are numbered 1 and 2; Grace's and
Linus's orders are numbered 1. An aggregate collapses rows into groups. A
window expression calculates a value across related rows while retaining
the row-level result. Consult the window chapter for the implementation's
fixed-frame behavior; explicit frame syntax is unavailable.

\section{Workshop 4: transactions and targeted changes}
\subsection{Tasks}
Reduce keyboard stock by one. Create a savepoint. Reduce it by two more.
Undo only the second change, verify the stock, then undo the entire
transaction. Explain why the final value must equal its initial value.
Do this in one shell session or one script.

\subsection{Solution and checkpoints}
\begin{lstlisting}[language=SQL]
BEGIN;
UPDATE products SET stock = stock - 1 WHERE product_id = 101;
SAVEPOINT second_attempt;
UPDATE products SET stock = stock - 2 WHERE product_id = 101;
SELECT stock FROM products WHERE product_id = 101;
ROLLBACK TO second_attempt;
SELECT stock FROM products WHERE product_id = 101;
ROLLBACK;
SELECT stock FROM products WHERE product_id = 101;
\end{lstlisting}
The displayed values are 9, 11 and 12. The first rollback restores the
savepoint's state; it does not end the transaction. The final rollback
restores the transaction's initial state. Two independent \code{--query}
commands do not share this state, even if both connect to the same port.

As a second exercise, first \code{SELECT} the rows matching a proposed
\code{UPDATE} or \code{DELETE} condition. Compare their identifiers with
your intention before running the change. This habit catches an omitted
predicate more reliably than checking a large affected-row count afterwards.

\section{Workshop 5: application input and result decoding}
\subsection{Tasks}
Write a Python function that fetches a customer by identifier using a bound
parameter. Close both cursor and connection. Explain the return type of the
identifier. Extend the function to return \code{None} when the customer does
not exist. Keep connection secrets outside source code.

\subsection{Solution}
\begin{lstlisting}[language=Python]
import minisql

def customer_by_id(dsn, customer_id):
    connection = minisql.connect(dsn, autocommit=True)
    try:
        with connection.cursor() as cursor:
            cursor.execute(
                "SELECT customer_id, name FROM customers "
                "WHERE customer_id = ?",
                (customer_id,),
            )
            row = cursor.fetchone()
            if row is None:
                return None
            return {"customer_id": int(row[0]), "name": row[1]}
    finally:
        connection.close()
\end{lstlisting}
The connector returns ordinary non-null result cells as strings, so the
application converts the known integer identifier explicitly. A bound
parameter represents a value; it does not replace a table or column name.
The comma in \code{(customer\_id,)} creates a one-element Python tuple.

The connection context manager in this connector controls a transaction
boundary and does not close the connection. The explicit \code{finally}
block is intentional. With \code{autocommit=True}, this single read does not
leave a lazy explicit transaction open. For a multi-step write operation,
use the transaction pattern in the connector chapter instead.

\section{Workshop 6: investigate an index}
\subsection{Tasks}
Explain a query that looks up a product by primary key. Explain a query that
filters a non-key column. On a separate larger fixture, create a suitable
explicit index, run \code{ANALYZE}, then compare \code{EXPLAIN ANALYZE}
before and after. Record both plan shape and observed work.

\subsection{Procedure}
\begin{lstlisting}[language=SQL]
EXPLAIN SELECT product_id, name
FROM products WHERE product_id = 103;

EXPLAIN SELECT product_id, name
FROM products WHERE stock = 100;

CREATE INDEX idx_product_stock ON products(stock);
ANALYZE products;

EXPLAIN ANALYZE SELECT product_id, name
FROM products WHERE stock = 100;
\end{lstlisting}
The acceptance condition is not a particular millisecond value or a promise
that a five-row table must use an index. You should be able to identify the
selected access path, describe the predicate, and explain the trade-off
between reading the table and using a key followed by heap access. Repeat on
data of realistic size and distribution before drawing a performance
conclusion. The performance chapter covers covering, partial and functional
indexes, the full-equality requirement for composite index access, and
statistics that affect these choices.

\section{Workshop 7: a backup is useful only if you can restore it}
\subsection{Tasks}
Use a disposable teaching database. Stop its server gracefully. Create a
backup into a new location using the exact command in the backup chapter.
Restore into a different new location. Run the offline checker on the
restored database, start it on a different port, and compare table counts
and representative joined results with the source.

\subsection{Acceptance conditions}
Record the source directory, backup path, restored directory, source
revision, storage format, page size and tool exit statuses. If encryption
is enabled, demonstrate that the separately protected matching key is
available. The restored server must produce counts 4, 5, 4 and 6 for the
teaching tables, plus Ada 6 and Grace 3 for the grouped-units query. Make a
small reversible write on the restored copy to prove that it is usable as
the intended writable recovery target.

A successful command that merely creates an archive is not sufficient.
An untested key, missing WAL range or schema-mismatched replication base can
remain invisible until recovery is needed. The goal is to prove a full
procedure without changing the original database.

\section{Workshop 8: explain a current implementation limit}
Choose one of the limitations below and write a short explanation containing
the intended SQL meaning, observed behavior at this revision, an example
that exposes the difference and a practical design response.
\begin{itemize}
\item Scaled-decimal coefficients, arithmetic and wire formatting.
\item A parent-key update that leaves a foreign-key child unchanged.
\item \code{CURRENT\_TIMESTAMP} returning epoch zero.
\item \code{NULLIF(5,NULL)} returning null in this implementation.
\item Automatic identities using the current maximum row value rather than
an independent durable sequence allocator.
\item Text \code{'NULL'} and SQL NULL being ambiguous in legacy protocol-v1
results, while peers negotiating results-v2 preserve the distinction.
\end{itemize}
The worked probe results in the validation appendix are the answer key.
An expert answer states the boundary precisely. For example, binary64 row
storage does not fix decimal arithmetic, and a successful foreign-key
definition does not prove that cascading actions execute. When a future
revision repairs a limitation, rerun the relevant probe and revise both
the warning and the expected result.

\section{A capstone project}
Build a small inventory application around the shop model. It should have
parameterized customer/product lookups, a transaction for an order and its
lines, immutable referenced keys, a documented money representation,
explicit handling of every database error, and a repeatable restore test.
Include a short operator guide with the server mode, credentials source,
resource limits, log location, backup retention policy and stop procedure.
Document the exact MiniSQL revision you used.

Assess the project by asking whether a second person can reproduce it from
the instructions, whether failed operations leave a verified state, and
whether the application can distinguish a known failure from an uncertain
commit after a broken connection. These checks combine language knowledge
with the operational judgement needed to use a database well.
\sourcecode{clients/python/minisql/dbapi.py}{Parameter binding, cursor behavior, lazy transactions and connection lifetime.}
\sourcecode{src/minisql/planner/optimizer.ml}{Costed access paths and the index eligibility rules used in the performance workshop.}
